En este informe se analizaron experimental y teóricamente el comportamiento de dos dispositivos cerrados de dos puertos (cuadripolos) con el propósito de determinar sus parámetros característicos. Utilizando un generador de señales y un osciloscopio, se midieron de forma individual las matrices de parámetros de impedancia ($Z$), admitancia ($Y$) y transmisión ($T$) a una frecuencia de excitación constante de $1\text{ kHz}$. A partir de dichas matrices, se dedujeron los componentes de los modelos circuitales equivalentes ($\Pi$ y $T$) propuestos por la cátedra para aproximar la topología interna desconocida de cada dispositivo.

A través de la teoría de conversión de parámetros, se validaron los resultados experimentales contrastando las matrices obtenidas de forma directa en el laboratorio con aquellas derivadas mediante transformaciones analíticas basadas en la matriz de impedancias. Asimismo, se estudiaron las propiedades de interconexión de redes combinadas bajo las topologías en cascada, serie y paralelo; para estas dos últimas, se verificó rigurosamente el cumplimiento previo de la condición de Brune para asegurar la validez del modelado matricial elemental. Adicionalmente, se evaluó el efecto de una carga externa sobre el cuadripolo mediante el empleo de funciones de red y se analizó la validez de los parámetros en función de la frecuencia en un rango de $100\text{ Hz}$ a $10\text{ kHz}$. En líneas generales, los resultados experimentales exhibieron una sólida concordancia con las simulaciones en LTspice y los cálculos teóricos, permitiendo convalidar los modelos matriciales lineales aplicados a redes de dos puertos.
